\documentclass[11pt]{article}
\usepackage[a4paper,margin=28mm]{geometry}
\usepackage{amsmath,amssymb,amsthm,mathtools}
\usepackage{booktabs,array,microtype}
\usepackage[dvipsnames]{xcolor}
\usepackage[hidelinks]{hyperref}
\usepackage{enumitem}
\setlist{nosep}
\newtheorem{theorem}{Theorem}[section]
\newtheorem{lemma}[theorem]{Lemma}
\newtheorem{proposition}[theorem]{Proposition}
\newtheorem{corollary}[theorem]{Corollary}
\theoremstyle{remark}
\newtheorem{remark}[theorem]{Remark}
\newcommand{\N}{\mathbb N}
\newcommand{\Q}{\mathbb Q}
\newcommand{\Z}{\mathbb Z}
\newcommand{\C}{\mathbb C}
\newcommand{\ph}{\varphi}
\DeclareMathOperator{\Res}{Res}
\numberwithin{equation}{section}

\title{\textbf{Coefficient Propagation and Conjugate Obstructions\\
for Trinomial Factors of Newman Polynomials}\\[3mm]
\large Local certificates, arithmetic norm divisors, and degree obstructions}
\author{Kende Daniel Rethelyi}
\date{28 July 2026}

\begin{document}
\maketitle

\begin{center}
\small
\textbf{Autonomous AI research notice.}
This manuscript was researched, critically revised, and generated
autonomously through an AI Researcher--Critic workflow.
\end{center}
\medskip

\begin{abstract}
The unfair $0$--$1$ polynomial conjecture asks whether monic polynomials with
nonnegative real coefficients must themselves be $0$--$1$ polynomials whenever
their product is a $0$--$1$ polynomial.  In a normalized nonnegative
factorization of a Newman polynomial, we study a prospective trinomial factor
$A_a(x)=1+a x^u+x^v$.  Coefficient propagation settles the cases in
which the two gaps are comparable by divisibility and forces every genuinely
fractional coefficient into $0<a<1/2$.  Algebraic integrality and the
golden-ratio annulus for zeros of Newman polynomials give explicit bounds on
all conjugates.  The complete norm polynomial is an integer divisor of the
product, yielding finite searches at fixed algebraic degree.  A unit-circle
crossing argument excludes quadratic coefficients when the smaller gap is
$2$.  Finally, an exact rational Sturm--Schur sieve and a parity-block
expansion of the forced cofactor recurrence exclude every cubic coefficient
in every remaining smaller-gap-$2$ case, with no product-degree cap.  An exact
quartic norm sieve and an expanded parity-block certificate also exclude every
quartic coefficient, giving under these factorization hypotheses the
algebraic-degree lower bound $[\Q(a):\Q]\ge5$.  The finite cubic catalog has
$126$ entries.  Exact algebraic signs settle $v=5,7,\ldots,21$, while rational
interval bounds for alternating binomial tails prove uniform positivity
followed by a negative coefficient for every odd $v\ge23$.  For quartics the
exact sieve retains $6{,}977$ minimal polynomials and $6{,}982$ target roots;
the corresponding negative coefficient occurs no later than degree
$110v+100$.  Independently, a unit-disk/Pringsheim argument proves eventual
negativity of the forced series for every real $a>0$.  A safe-lattice theorem
proves that the only possible coefficient-one escape cannot first occur before
degree $13v+4$ at an odd index or $14v+2$ at an even index.  These are
structural restrictions, not a proof of the full conjecture: the safe-lattice
theorem does not rule out a later escape, and the degree-$5$ and higher cases
remain open.
\end{abstract}

\section{Introduction}

A \emph{Newman polynomial} has coefficients in $\{0,1\}$ and constant term
$1$.  We consider the following open problem.

\medskip
\noindent\textbf{Unfair $0$--$1$ polynomial conjecture.}
If $A,B\in\mathbb R[x]$ are monic with nonnegative coefficients and $AB$ is a
$0$--$1$ polynomial, then both $A$ and $B$ are $0$--$1$ polynomials.
\medskip

The conjecture is equivalent to a factorization question for finitely
supported probability measures.  The arithmetic-progression case goes back
to Krasner and Ranulac~\cite{KR}.  Ghidelli treated the first difficult
trinomial $1+a x^2+x^5$ by recurrence and root analysis~\cite{Ghidelli}, and
Hare developed recursive symbolic methods for bounded-degree candidate
factors~\cite{Hare}.  The general conjecture remains open.

This paper isolates a uniform local coefficient mechanism and joins it to the
arithmetic theory of Newman divisors.  Throughout, the factor statements concern
normalized factorizations with nonnegative real cofactors.  The principal
conclusions are:
\begin{enumerate}
\item divisibility-comparable gaps force fairness;
\item a fractional middle coefficient is an irrational algebraic integer in
      $(0,1/2)$;
\item every conjugate satisfies explicit gap-dependent bounds;
\item the full norm polynomial is an effectively computable integer divisor;
\item smaller gap $2$ excludes algebraic degrees $2$, $3$, and $4$ without a
      product-degree restriction;
\item in the unrestricted smaller-gap-$2$ recurrence, a first coefficient-one
      escape is impossible below $13v+4$ at odd indices and $14v+2$ at even
      indices.
\end{enumerate}
The last assertion is an analytic theorem in the same normalized
nonnegative-factorization setting, accompanied by reproducible exact identity
and symbolic sign audits.  Its artifacts include source code, exact root
records, exact algebraic signs, rational interval margins, alternating-tail
bounds, and standalone checkers.
No priority claim is made for the general propagation or norm-divisor
mechanisms.

\section{Normalization and the recurrence}

\begin{lemma}[Endpoint normalization]\label{lem:normalization}
Let $A,B\in\mathbb R[x]$ be nonzero and monic with nonnegative coefficients,
and suppose $C=AB$ is a $0$--$1$ polynomial.  After removing powers of $x$
from $A$ and $B$, both constant coefficients are $1$, and every coefficient
of $A$ and $B$ belongs to $[0,1]$.
\end{lemma}
\begin{proof}
Let $a_0,b_0$ be the first nonzero coefficients after removing the powers of
$x$.  The first nonzero coefficient of $C$ is $a_0b_0=1$.  The coefficient in
degree $\deg B$ contains the nonnegative summand $a_0$, so $a_0\leq1$;
similarly $b_0\leq1$.  Hence $a_0=b_0=1$.  Multiplication by the constant
coefficient of the other factor makes each coefficient of $A$ or $B$ a
nonnegative summand of a coefficient of $C$, and hence at most $1$.
\end{proof}

If one normalized factor is already a $0$--$1$ polynomial, low-to-high
coefficient comparison makes the other factor integral; Lemma
\ref{lem:normalization} then makes it a $0$--$1$ polynomial as well.

For
\[
 A_a(x)=1+a x^u+x^v,\qquad 0<u<v,
\]
write $B(x)=\sum_{n\ge0}b_nx^n$ and $C(x)=\sum_{n\ge0}c_nx^n$, extending
$b_n=0$ outside the degree of $B$.  Then
\begin{equation}\label{eq:recurrence}
 c_n=b_n+a b_{n-u}+b_{n-v}\in\{0,1\},\qquad b_n\ge0,
\end{equation}
with $b_0=1$ and $b_j=0$ for $j<0$.  Whenever one summand in
\eqref{eq:recurrence} is positive, $c_n$ is forced to be $1$.

\section{Local propagation along arithmetic rays}

\begin{lemma}[Arithmetic-ray overflow]\label{lem:overflow}
Let $r\ge1$, $k\ge2$, and $0<a<1$.  Suppose a polynomial or formal series
begins
\[
 A(x)=1+a x^r+x^{kr}+O(x^{kr+1}),
\]
with no other nonzero coefficient through degree $kr$.  There is no series
$B=\sum_{j\ge0}b_jx^j$ with $b_0=1$, $b_j\ge0$, for which all coefficients of
$AB$ through degree $kr$ lie in $\{0,1\}$.
\end{lemma}
\begin{proof}
At degrees $r,2r,\ldots,(k-1)r$, positivity forces the product coefficient to
be $1$, and therefore
\[
 b_{jr}=1-a b_{(j-1)r}=1-a+a^2-\cdots+(-a)^j
       =\frac{1-(-a)^{j+1}}{1+a}>0.
\]
At degree $kr$, the terminal unit term collides with this positive tail:
$[x^{kr}](AB)=b_{kr}+a b_{(k-1)r}+b_0>1$.
\end{proof}

\begin{theorem}[Divisible gaps]\label{thm:divisible}
Suppose $A_a=1+a x^u+x^v$ occurs in a normalized nonnegative factorization of
a $0$--$1$ polynomial.  If one of $u$ and $v-u$ divides the other, then
$a\in\{0,1\}$ and both factors are $0$--$1$ polynomials.
\end{theorem}
\begin{proof}
If $u\mid v$, apply Lemma~\ref{lem:overflow} with $r=u$.  If $(v-u)\mid v$,
reverse the factorization using $P^*(x)=x^{\deg P}P(x^{-1})$; the reversed
trinomial is $1+a x^{v-u}+x^v$.  These two alternatives are equivalent to
comparability of $u$ and $v-u$ by divisibility.  Thus $0<a<1$ is impossible,
and the observation after Lemma~\ref{lem:normalization} proves the final claim.
\end{proof}

\begin{theorem}[Universal half-bias obstruction]\label{thm:half}
If $A_a$ occurs in a normalized nonnegative factorization of a $0$--$1$
polynomial and $0<a<1$, then $a\le1/2$.
\end{theorem}
\begin{proof}
After reversal assume $u\le v-u$.  Equality is covered by Theorem
\ref{thm:divisible}, so $2u<v$.  Recurrence~\eqref{eq:recurrence} gives
$b_u=1-a$ and $b_{2u}=1-a+a^2$.  At degree $v$, the unit contribution from
$x^vb_0$ forces $b_v=b_{v-u}=0$; the next position gives $b_{v+u}=a$.
Consequently
\[
[x^{v+2u}](A_aB)=b_{v+2u}+a b_{v+u}+b_{2u}
=b_{v+2u}+1-a+2a^2.
\]
For $a>1/2$ this exceeds $1$, a contradiction.
\end{proof}

\begin{proposition}[Forced-ray skeleton]\label{prop:skeleton}
Assume the hypotheses of Theorem~\ref{thm:half}, with $2u<v$.  By
Theorem~\ref{thm:divisible}, the feasible fractional case has $u\nmid v$.
Thus write $v=qu+r$, where $q\ge2$ and $0<r<u$.  Put
\[
 S_j(a)=\sum_{k=0}^j(-a)^k,\qquad
 T_j(a)=\sum_{k=0}^{j-1}(k+1)(-a)^k.
\]
Every feasible cofactor satisfies
\[
 b_{ju}=S_j(a)\ (0\le j\le q),\quad
 b_{r+ju}=0\ (0\le j\le q-1),\quad
 b_{v+ju}=aT_j(a)\ (1\le j\le q),
\]
and
\[
 T_j(a)=\frac{1-(j+1)(-a)^j+j(-a)^{j+1}}{(1+a)^2}.
\]
For $0<a\le1/2$, every $T_j(a)$ is nonnegative; for $0<a<1/2$ it is
positive.  Thus this paired ray gives no stronger upper bound before the next
interaction at degree $2v$.
\end{proposition}
\begin{proof}
The first identity is the recurrence below degree $v$.  At degree $v$ one has
$b_{v-u}=0$.  If a zero $b_{r+ju}=0$ in this residue class had
$b_{r+(j-1)u}>0$, then the recurrence below degree $v$ would force
$a b_{r+(j-1)u}=1$, hence $b_{r+(j-1)u}=1/a>1$, contrary to
Lemma~\ref{lem:normalization}.  Backward induction gives the displayed zero
chain.  Put $d_j=b_{v+ju}$ and $d_0=0$.  Since $S_j(a)>0$,
the product coefficient at $v+ju$ is $1$, whence
$d_j=a(S_{j-1}(a)-d_{j-1})$.  Induction yields $d_j=aT_j(a)$, and
differentiating a finite geometric sum gives the closed form.  For
$0<a<1/2$, the magnitudes $(k+1)a^k$ decrease strictly, so the finite
alternating sum is positive.  At $a=1/2$ they are nonincreasing and the only
zero case is $T_2(1/2)=0$; hence nonnegativity holds on the closed interval.
\end{proof}

\section{Algebraic integrality and the norm divisor}

\begin{theorem}[Algebraic integrality]\label{thm:integral}
Let $C\in\Z[x]$ be monic and $C=AB$ with $A,B\in\mathbb R[x]$ monic.  Every
coefficient of $A$ and $B$ is an algebraic integer.  In particular, a rational
coefficient is an integer.
\end{theorem}
\begin{proof}
Every root of either factor is a root of the monic integer polynomial $C$ and
is therefore an algebraic integer.  Factor coefficients are elementary
symmetric polynomials in these roots.  A rational algebraic integer is an
ordinary integer.
\end{proof}

Thus a genuinely fractional coefficient satisfies
\begin{equation}\label{eq:fractional}
 0<a<\frac12,\qquad a\text{ is an irrational algebraic integer}.
\end{equation}

\begin{theorem}[Conjugate and norm sieve]\label{thm:norm}
Suppose $A_a(x)=1+a x^u+x^v$ occurs in a normalized nonnegative factorization
$C=A_aB$, where $C$ is a Newman polynomial.  Let $a$ have degree $d$, let
$m_a(t)$ be its minimal polynomial, and put $\ph=(1+\sqrt5)/2$.  For every
conjugate $\alpha$ of $a$:
\begin{enumerate}[label=(\alph*)]
\item $A_\alpha(x)=1+\alpha x^u+x^v$ divides $C$ in $\overline{\Q}[x]$;
\item
\begin{equation}\label{eq:Rphi}
 |\alpha|\le R_\ph(u,v):=
 \min\{\ph^u+\ph^{u-v},\ \ph^{v-u}+\ph^{-u}\};
\end{equation}
\item if $\alpha$ is real, then $\alpha>-\kappa(u,v)$, where
\begin{equation}\label{eq:kappa}
 \kappa(u,v)=
 \left(\frac{v-u}{u}\right)^{u/v}
 +\left(\frac{u}{v-u}\right)^{(v-u)/v}.
\end{equation}
\end{enumerate}
Furthermore,
\begin{equation}\label{eq:norm}
 H_a(x):=\prod_\alpha(1+\alpha x^u+x^v)
 =\Res_t\bigl(m_a(t),1+t x^u+x^v\bigr)
\end{equation}
belongs to $\Z[x]$, divides $C$ in $\Z[x]$, and has degree $dv$.
Consequently $\deg C\ge dv$.
\end{theorem}
\begin{proof}
Division of $C$ by the monic polynomial $A_a$ takes place in $\Q(a)[x]$.
Applying an embedding $\Q(a)\hookrightarrow\C$ to $C=A_aB$ fixes $C$ and
proves (a).

Odlyzko and Poonen's zero bound gives
\begin{equation}\label{eq:annulus}
 \ph^{-1}<|z|<\ph
\end{equation}
for every zero of a Newman polynomial~\cite{OP}.  On $|z|=\ph^{-1}$, the
strict inequality $|\alpha|\ph^{-u}>1+\ph^{-v}$ would make $\alpha z^u$
dominate $1+z^v$; Rouch\'e's theorem would place $u$ zeros of $A_\alpha$ in
the forbidden inner disk.  Hence $|\alpha|\le\ph^u+\ph^{u-v}$.  On
$|z|=\ph$, the analogous strict inequality would make $A_\alpha$ have only
$u$ zeros in that circle, whereas all $v$ of its zeros must lie there by (a)
and~\eqref{eq:annulus}.  This gives the other bound in~\eqref{eq:Rphi}.

For (c), write $\alpha=-s$ with $s>0$.  The function
$f(x)=1-sx^u+x^v$ on $x>0$ first acquires a positive zero at tangency.
Solving $f(x)=f'(x)=0$ gives $s=\kappa(u,v)$.  Thus for
$s\ge\kappa(u,v)$, $A_{-s}$ has a positive real zero.  Since $A_{-s}$ divides
$C$, that zero would also be a zero of $C$, impossible because $C$ has
nonnegative coefficients and $C(0)=1$.  Notice that no coefficient-sign claim
about $A_{-s}$ is used.

Distinct conjugate trinomials are coprime: their difference is a nonzero
multiple of $x^u$, and neither vanishes at $0$.  Their product is Galois-fixed
and has algebraic-integer coefficients, hence lies in $\Z[x]$.  By (a) and
coprimality it divides $C$ in $\overline{\Q}[x]$, hence in $\Q[x]$; Gauss's
lemma gives divisibility in $\Z[x]$.  Formula~\eqref{eq:norm} and the degree
claim follow directly.
\end{proof}

\section{Effective finiteness at fixed algebraic degree}

\begin{theorem}[Effective finite enumeration]\label{thm:finite}
Fix $u,v$ and $d$.  There are only finitely many algebraic integers $a$ of
degree $d$ that can occur in~\eqref{eq:fractional}.  More explicitly, if
\[
 m_a(t)=t^d+s_{d-1}t^{d-1}+\cdots+s_0\in\Z[t],
\]
and $R=R_\ph(u,v)$, then
\begin{equation}\label{eq:boxes}
 |s_{d-k}|\le {d\choose k}R^k,\qquad 1\le k\le d.
\end{equation}
Every candidate can therefore be enumerated in finite integer boxes and
tested against the necessary root and conjugate conditions.  For any fixed
product-degree cap $D$, existence of a Newman multiple is a finite integer
feasibility problem for the quotient coefficients.  If $\deg C\le D$, it is
enough to consider $d\le\lfloor D/v\rfloor$.
\end{theorem}
\begin{proof}
Up to sign, $s_{d-k}$ is the $k$th elementary symmetric function of the $d$
conjugates.  It has exactly ${d\choose k}$ summands, each of modulus at most
$R^k$, which proves~\eqref{eq:boxes}.  The coefficient boxes are finite.  For
a fixed candidate $H_a$ and cap $D$, write $C=H_aQ$ with
$\deg Q\le D-dv$.  Since $H_a(0)=C(0)=1$, one has $q_0=1$, and low-to-high
coefficient recursion
\[
 q_n=c_n-\sum_{i=1}^n h_iq_{n-i}
\]
shows inductively that every $q_n$ is integral.  Equating the finitely many
coefficients and requiring each coefficient of $C$ to lie in $\{0,1\}$ then
gives the stated finite integer system.
The degree restriction follows from $dv\le\deg C$.
\end{proof}

\section{The smaller-gap-two quadratic case}

\begin{lemma}\label{lem:kappatwo}
For every integer $v\ge5$ one has $\kappa(2,v)<2$.
\end{lemma}
\begin{proof}
Put $p=2/v$ and $r=p/(1-p)=2/(v-2)$.  Then $0<r\le2/3$ and
\[
 \kappa(2,v)=r^{-p}+r^{1-p}=(1+r)r^{-r/(1+r)}=:K(r).
\]
Direct differentiation gives
\[
 \frac{d}{dr}\log K(r)=-\frac{\log r}{(1+r)^2}>0
 \qquad(0<r<1).
\]
Hence $K(r)\le K(2/3)=\frac53(3/2)^{2/5}<2$.  The final inequality is
equivalent, after taking fifth powers, to
$9/4<(6/5)^5$, i.e. $28125<31104$.
\end{proof}

\begin{proposition}[Unique coarse quadratic survivor]\label{prop:quadratic}
Suppose $\min\{u,v-u\}=2$.  If a genuinely fractional coefficient $a$ is
quadratic, then necessarily $a=(3-\sqrt5)/2$.
\end{proposition}
\begin{proof}
Reverse if necessary, so $u=2\le v-u$.  Theorem~\ref{thm:divisible} removes
even $v$ and equality, so $v\ge5$ is odd.  Let $a'$ be the other conjugate.
If the integer norm $aa'$ were negative, then $a'\le-1/a<-2$.  On the other
hand Theorem~\ref{thm:norm}(c) and Lemma~\ref{lem:kappatwo} give
$a'>-\kappa(2,v)>-2$, a contradiction.  Thus $a'>0$.  The modulus bound gives
\[
 a'\le\ph^2+\ph^{2-v}\le\ph^2+\ph^{-3}
     =\frac{3\sqrt5-1}{2}<2.855.
\]
Writing the minimal polynomial as $t^2-Tt+N$, one obtains
$0<N=aa'<2$ and $0<T=a+a'<4$.  Hence $N=1$ and $T\le3$.  The only monic
quadratic with a positive root in $(0,1/2)$ is $t^2-3t+1$.
\end{proof}

\begin{lemma}[Boundary-root perturbation]\label{lem:boundary}
Let $v\ge5$ be odd and $0<t<2\sin(\pi/(v-1))$.  Then
\[
 g_t(z)=1+z^2+t z^v
\]
has exactly two zeros in the open unit disk.
\end{lemma}
\begin{proof}
At $t=0$ the two roots are $\pm i$.  The implicit-function theorem gives an
analytic branch from $i$ with
\[
 z(t)=i-\frac{i^{v-1}}2t-\frac{2v-1}{8}it^2+O(t^3),
 \qquad
 |z(t)|^2=1-\frac{v-1}{2}t^2+O(t^3).
\]
Thus this branch and its conjugate branch from $-i$ lie inside the unit circle
for all sufficiently small $t>0$.  They are the only bounded branches: if
$t_n\to0$ and zeros $z_n$ remain bounded, every limit point solves
$1+z^2=0$; the two simple implicit branches account for all such zeros.
Hence, for sufficiently small positive $t$, exactly two zeros lie in the unit
disk.

It remains to rule out boundary crossings.  If $g_t(e^{i\theta})=0$, then
multiplication by $e^{-i\theta}$ gives
\[
 2\cos\theta+t e^{i(v-1)\theta}=0.
\]
Therefore $(v-1)\theta=k\pi$ and necessarily
\[
 t=-2(-1)^k\cos\!\left(\frac{k\pi}{v-1}\right)>0.
\]
Because $v-1$ is even, the zero cosine occurs at $\theta=\pi/2$ modulo
$\pi$, and the nearest nonzero point of the grid has absolute cosine
$\sin(\pi/(v-1))$.  Thus no unit-circle zero exists for
$0<t<2\sin(\pi/(v-1))$.  The argument principle now makes the number of
zeros in the unit disk constant throughout this interval, so it remains two.
\end{proof}

\begin{theorem}[Quadratic exclusion]\label{thm:noquadratic}
If $1+a x^u+x^v$ is a genuinely fractional factor in a normalized
nonnegative factorization of a Newman polynomial and
$\min\{u,v-u\}=2$, then $[\Q(a):\Q]\ge3$.
\end{theorem}
\begin{proof}
Assume $a$ is quadratic.  Proposition~\ref{prop:quadratic} gives
$a=(3-\sqrt5)/2$ and conjugate $a'=(3+\sqrt5)/2=\ph^2$.  Reverse so $u=2$
and $v\ge5$ is odd.  Under $z=\ph x$,
\[
 A_{a'}(z/\ph)=1+z^2+\ph^{-v}z^v.
\]
Moreover $\ph^{-v}<2\sin(\pi/(v-1))$: indeed
$2\sin(\pi/(v-1))\ge4/(v-1)$, while $\ph^v>F_v\ge v-1$ for $v\ge5$.
Lemma~\ref{lem:boundary} yields a zero $|z|<1$, hence a zero
$|x|<\ph^{-1}$ of $A_{a'}$.  Theorem~\ref{thm:norm}(a) makes this a zero of
$C$, contradicting~\eqref{eq:annulus}.
\end{proof}

\section{Global cubic exclusion for the smaller gap two}
\label{sec:globalcubic}

This section replaces a bounded quotient search by a forced-recurrence
certificate.  The argument has two exact computational layers.  First, the
uniform norm sieve produces the same $126$ possible cubic minimal polynomials
for every odd $v\ge5$.  Second, for the distinguished root $a\in(0,1/2)$ of
each polynomial, a parity-block expansion proves that the cofactor recurrence
stays strictly positive, apart from its forced initial zeros, until a specified
coefficient becomes negative.  The obstruction index is at most $27v+8$.

\begin{remark}[Why the nonnegative cofactor is essential]\label{rem:scopeglobal}
Let $a=2\cos(4\pi/9)$, the root in $(0,1/2)$ of the irreducible polynomial
$t^3-3t+1$.  Exact multiplication, using $a^3=3a-1$, gives
\[
 1+x^{18}=(1+a x^2+x^4)B(x),
\]
where
\begin{align*}
B(x)={}&x^{14}-a x^{12}+(a^2-1)x^{10}+(1-a)x^8\\
      &+(1-a)x^6+(a^2-1)x^4-a x^2+1.
\end{align*}
Thus cubic ordinary Newman divisibility does occur.  The negative coefficients
of $B$ show why the normalized nonnegative-factorization hypothesis cannot be
discarded.
\end{remark}

\begin{lemma}[Uniform cubic catalog]\label{lem:cubiccatalog}
If $A_a(x)=1+a x^2+x^v$ occurs in a normalized nonnegative factorization,
$v\ge5$ is odd, and $[\Q(a):\Q]=3$, then the minimal polynomial of $a$ is one
of the $126$ polynomials in the accompanying exact catalog.
\end{lemma}
\begin{proof}
Write $m_a(t)=t^3+s_2t^2+s_1t+s_0$.  For every odd $v\ge5$,
Theorem~\ref{thm:norm} gives
\[
 |\alpha|<\ph^2+\ph^{2-v}\le \ph^2+\ph^{-3}<\frac{143}{50}
\]
for every conjugate $\alpha$ of $a$.  Therefore
\[
 |s_2|\le8,\qquad |s_1|\le24,\qquad |s_0|\le23.
\]
If $\alpha,\beta$ are the other two conjugates, then
$|s_0|=|a\alpha\beta|$ is a positive integer.  Hence the target root also has
the uniform lower bound
\[
 a>\left(\frac{50}{143}\right)^2=\frac{2500}{20449}.
\]
The real-conjugate estimate also gives $\alpha>-\kappa(2,v)>-2$ for every
real conjugate.  The exact generator enumerates all $38318$ triples in the
boxes with nonzero constant coefficient, removes reducible cubics by the
integer-root test, counts roots in $(0,1/2)$ and $(-3,-2]$ by rational Sturm
sequences, and applies rational Schur recursion for the disk bound.  Exactly
$126$ irreducible cubics remain.  The independent checker repeats this
enumeration without importing the generator.
\end{proof}

For the recurrence certificate, define the formal series
\begin{equation}\label{eq:forcedseries}
 Q_{a,v}(x):=\frac{1+x^v}{(1-x^2)(1+a x^2+x^v)}
             =\sum_{n\ge0}q_nx^n.
\end{equation}
The eventual sign change in this series is not special to algebraic degree
three.
\begin{lemma}[Analytic sign change]\label{lem:pringsheim}
For every $a>0$ and every odd $v\ge5$, the series $Q_{a,v}$ has a negative
coefficient.
\end{lemma}
\begin{proof}
Put $A(x)=1+ax^2+x^v$.  Since $A$ is monic with constant coefficient $1$,
the product of the moduli of its roots is $1$.  If $A$ had no root in the open
unit disk, all its roots would therefore lie on the unit circle.  Reality of
the coefficients would then make the root multiset invariant under inversion,
so the monic reciprocal $x^vA(x^{-1})=x^v+ax^{v-2}+1$ would equal $A$.
This is impossible for $a>0$ and $v\ge5$.  Thus $A$ has a root of modulus
strictly less than $1$.

No zero of $A$ cancels against $1+x^v$: a common zero would give $ax^2=0$.
Nor does $A$ vanish at $1$ or $-1$.  Hence the radius of convergence of
$Q_{a,v}$ is some $R<1$.  If every $q_n$ were nonnegative, Pringsheim's
theorem would make the positive real point $R$ a singularity.  But the
displayed rational function has no singularity on $(0,1)$, since both
$1-x^2$ and $1+ax^2+x^v$ are positive there.  This contradiction proves the
claim.
\end{proof}
\begin{remark}[General nonpalindromic gaps]\label{rem:generalpringsheim}
The same proof works with $2$ replaced by any $u$ with $0<u<v$ and $2u\ne v$:
the formal series
\[
 \frac{1+x^v}{(1-x^u)(1+a x^u+x^v)}
\]
has a negative coefficient for every $a>0$.  Indeed, if all roots of the
trinomial lay on the unit circle, reciprocity would force the middle exponents
$u$ and $v-u$ to agree.  This analytic statement alone does not prove that an
arbitrary finite cofactor follows the displayed forced-ray series up to its
first negative term.
\end{remark}

The numerator is not an arbitrary guess: if the cofactor coefficients remain
nonnegative, positivity forces the product coefficients to be $1$ on the even
ray and on the odd ray beginning at $v$.

\begin{lemma}[Forced-series obstruction]\label{lem:forcedseries}
Let $0<a<1$, let $v\ge5$ be odd, and suppose there is an index $N>v$ such
that the coefficients of~\eqref{eq:forcedseries} satisfy
\[
 0\le q_n<1\qquad(0<n<N),
\]
and $q_N<0$.  Then $1+a x^2+x^v$ has no finite monic nonnegative cofactor
whose product is a Newman polynomial.
\end{lemma}
\begin{proof}
Extend a prospective cofactor $B=\sum b_nx^n$ by zero outside its degree.
Endpoint normalization gives $0\le b_n\le1$ for every $n$.
The coefficient at degree $v$ contains $b_0=1$, so it forces
$b_v=b_{v-2}=0$.  If $b_j=0$ at an odd index $j<v$ but
$b_{j-2}>0$, then the coefficient at degree $j$ would force
$ab_{j-2}=1$, and hence $b_{j-2}=1/a>1$, a contradiction.  Backward
propagation along the odd ray therefore gives
$b_1=b_3=\cdots=b_{v-2}=0$.  Along the even ray below $v$, positivity forces
the product coefficient to be $1$ and hence gives the coefficients of
$Q_{a,v}$.

For $n\ge v$, the defining recurrence for $Q_{a,v}$ gives
\[
 a q_{n-2}+q_{n-v}=1-q_n.
\]
This forcing term is strictly positive for $n<N$ by hypothesis and is greater
than $1$ at $n=N$.  Thus the product coefficient is forced to be $1$, and
induction gives $b_n=q_n$ through
$\min\{N,\deg B\}$.  If $\deg B\ge N$, this yields the contradiction
$b_N=q_N<0$.  If $0<\deg B<N$, then the leading coefficient $b_{\deg B}=1$.
But $b_{\deg B}=q_{\deg B}<1$, again contradicting monicity.  Finally $B=1$
would leave the fractional coefficient $a$ in the product.
\end{proof}

\begin{proposition}[The coefficient-one escape]\label{prop:oneescape}
Fix $0<a<1$ and odd $v\ge5$.  Suppose that, through degree $n-1\ge v$,
a prospective nonnegative cofactor agrees with the forced series
$Q_{a,v}=\sum q_jx^j$.  Put
\[
 f_n:=a q_{n-2}+q_{n-v}.
\]
Then $f_n>1$ is an immediate contradiction; if $0<f_n\le1$, the Newman
product coefficient at degree $n$ is forced to be $1$ and the next cofactor
coefficient is forced to be $q_n=1-f_n$; and if $f_n=0$, the first deviation
from the forced series is possible.  In the last case $q_n=1$, and, while the
preceding coefficients are nonnegative,
\[
 f_n=0\quad\Longleftrightarrow\quad q_{n-2}=q_{n-v}=0.
\]
No such simultaneous predecessor-zero event occurs for $v<n\le2v$.
\end{proposition}
\begin{proof}
At every degree $n\ge v$, the target ray series has coefficient $1$, so
$q_n=1-f_n$.  If the actual product coefficient is $c_n\in\{0,1\}$, then
nonnegativity of the next cofactor coefficient gives
$b_n=c_n-f_n\ge0$.  Thus $f_n>0$ forces $c_n=1$, and gives the first two
alternatives; when $f_n=0$, either $c_n=1$ and $b_n=q_n=1$, or $c_n=0$ and
$b_n=0$, which is the first possible branch away from $Q_{a,v}$.  Because
$a>0$, a zero sum of the two nonnegative predecessors is equivalent to their
simultaneous vanishing.

For $v<n\le2v$, if $n$ is even then $q_{n-v}$ lies on the initial odd zero
ray whereas $q_{n-2}=1-a+a^2-\cdots$ is strictly positive for $0<a<1$.
If $n$ is odd, then $q_{n-v}$ is such a positive alternating geometric sum.
Hence $f_n>0$ throughout this range.
\end{proof}

\begin{lemma}[Parity rationalization]\label{lem:parityrationalization}
Put $z=x^2$ and define $r_m=r_m(t)$ by
\[
 R_{t,v}(z):=\sum_{m\ge0}r_mz^m
 =\frac{1}{(1-z)((1+tz)^2-z^v)},
 \qquad r_m=0\quad(m<0).
\]
Then the forced-series coefficients, regarded as polynomials in $t$, satisfy
\begin{align*}
 q_{2m}(t)&=r_m+t r_{m-1}-r_{m-v},\\
 q_{v+2m}(t)&=t r_{m-1}.
\end{align*}
Moreover
\[
 r_m+2t r_{m-1}+t^2r_{m-2}-r_{m-v}=1
 \qquad(m\ge0).
\]
\end{lemma}
\begin{proof}
Rationalizing with respect to the odd power $x^v$ gives
\begin{align*}
 Q_{t,v}(x)
 &=\frac{(1+x^v)(1+tx^2-x^v)}
 {(1-x^2)((1+tx^2)^2-x^{2v})}\\
 &=\frac{1+tz-z^v+t x^v z}
 {(1-z)((1+tz)^2-z^v)}.
\end{align*}
Separating the even part and the $x^v$ times an even series gives the two
coefficient identities.  The recurrence for $r_m$ is the coefficient of
$z^m$ in
$((1+tz)^2-z^v)R_{t,v}(z)=(1-z)^{-1}$.
\end{proof}

For $k,m\ge0$ put
\[
 P_{k,m}(t):=\sum_{\ell=0}^{m}(-1)^\ell
              {k+\ell\choose\ell}t^\ell,
 \qquad P_{k,\infty}(t)=(1+t)^{-k-1}.
\]
These are the Taylor partial sums of $(1+t)^{-k-1}$.  The next theorem
turns the parity rationalization into an infinite forbidden region for the
two-predecessor zero pair.

\begin{theorem}[Predecessor safe lattice]\label{thm:safelattice}
Let $0<t<1/2$, let $v\ge5$ be odd, and let $n>2v$.  If
\[
 q_{n-2}(t)=q_{n-v}(t)=0,
\]
then
\[
 n\ge5v+4\quad\hbox{when $n$ is odd},\qquad
 n\ge6v+2\quad\hbox{when $n$ is even}.
\]
Thus a first coefficient-one escape in Proposition~\ref{prop:oneescape}
obeys the same lower bounds.
\end{theorem}
\begin{proof}
Write $E_m=q_{2m}$.  Lemma~\ref{lem:parityrationalization} and its recurrence
give
\begin{equation}\label{eq:Er}
 E_m=1-t r_{m-1}-t^2r_{m-2}.
\end{equation}
If $n$ is odd, put $j=(n-v)/2-2$; if $n$ is even, put
$j=n/2-v-1$.  The predecessor pairs become, respectively,
\begin{equation}\label{eq:predecessorpairs}
 (q_{n-2},q_{n-v})=(t r_j,E_{j+2}),\qquad
 (q_{n-v},q_{n-2})=(t r_j,E_{j+v}).
\end{equation}
It is therefore enough to exclude a common zero in either pair for
$0\le j<2v$.

Expanding the defining rational function for $R_{t,v}$ geometrically in
$z^v/(1+tz)^2$ gives, for $0\le j<2v$,
\begin{equation}\label{eq:firsttworblocks}
 r_j=P_{1,j}+\boldsymbol{1}_{j\ge v}P_{3,j-v}.
\end{equation}
The terms $(\ell+1)t^\ell$ decrease strictly when $0<t<1/2$, so every
$P_{1,j}$ is positive.  Taylor's theorem also gives
$P_{k,2h}>P_{k,\infty}>0$.  Hence $r_j>0$ unless
$j=v+s$ with $s$ odd.  Odd $P_{3,s}$ increase with $s$, while for $s\ge7$
\[
 r_{v+s}> (1+t)^{-2}+P_{3,7}(t)>0.
\]
For completeness, the last function is decreasing on $(0,1/2)$: after
differentiation its negative numerator is twice
\[
 420t^9+1008t^8+644t^7+14t^6
 +(10t^4-12t^5)+(6t^2-8t^3)+(3-4t),
\]
which is positive term by term.  Its value at $t=1/2$ is $37/144$.
Consequently only $s\in\{1,3,5\}$ can remain.

First consider the odd pair in~\eqref{eq:predecessorpairs}.  At a zero of
$r_{v+s}$, equation~\eqref{eq:Er} gives
\[
 E_{v+s+2}=1-t r_{v+s+1}.
\]
Here $P_{1,v+s+1}<(1+t)^{-2}$, and the even partial sums $P_{3,2h}$
decrease from $P_{3,2}$ for $h\ge1$.  Thus
\[
 r_{v+s+1}<(1+t)^{-2}+P_{3,2}(t).
\]
The resulting lower bound
\[
 1-t\bigl((1+t)^{-2}+1-4t+10t^2\bigr)
\]
is positive on $(0,1/2)$.  Indeed, after multiplication by $(1+t)^2$ its
complement is
$10t^5+16t^4+3t^3-3t^2$, which is at most $15/16$: use
$10t^5\le5t^4$, $t^4\le t^2/4$, and $t^3\le t^2/2$.
Therefore the odd pair never vanishes for $j<2v$.

It remains to treat the even pair for the three exceptional offsets.  We
will use the exact identity
\begin{equation}\label{eq:evenexception}
 E_{2v+s}=P_{0,2v+s}
 -t\bigl(P_{4,s-1}+P_{2,v+s-1}\bigr),
 \qquad s\in\{1,3,5\},
\end{equation}
obtained from the same geometric expansion.  The odd partial sums
$P_{2,m}$ increase for odd $m\ge5$ and are positive: $P_{2,5}$ is decreasing
on $(0,1/2)$ because
\[
 P_{2,5}'=-3\{t^2(35t^2-20t+3)+(7t^2-4t+1)\}<0,
\]
and $P_{2,5}(1/2)=1/32$.

For $s=1$, a zero of $r_{v+1}$ satisfies
$P_{1,v+1}=4t-1$.  The even $P_{1,2h}$ decrease, so
$P_{1,6}-(4t-1)\ge0$.  This difference decreases on $(0,1/2)$ and at
$t=383/1000$ equals
\[
 -\frac{2424992339692617}{10^{18}}<0.
\]
Hence $t<383/1000$.  Since $P_{0,2v+1}>1-t$ and
$P_{2,v}<(1+t)^{-3}$, equation~\eqref{eq:evenexception} gives
\[
 E_{2v+1}>1-2t-\frac{t}{(1+t)^3}>0;
\]
the last expression decreases and at $t=383/1000$ is already
\[
 \frac{117994119779}{1322624443500}>0.
\]

For $s=3$, the function $(1+t)^{-2}+P_{3,3}(t)$ is decreasing.  At
$t=433/1000$ its value is
\[
 \frac{637744741107}{102674450000000}>0.
\]
Thus a zero of $r_{v+3}$ has $t>433/1000$.  Since
$P_{0,2v+3}<(1+t)^{-1}$ and the
$P_2$ term in~\eqref{eq:evenexception} is positive,
\[
 E_{2v+3}<(1+t)^{-1}-tP_{4,2}(t)<0.
\]
The right side decreases beyond $433/1000$ and at that endpoint is
\[
 -\frac{4430649363}{286600000000}.
\]

For $s=5$, direct differentiation shows that
$(1+t)^{-2}+P_{3,5}(t)$ is decreasing; its derivative has negative
numerator twice
\[
 140t^7+350t^6+240t^5+10t^4+(6t^2-8t^3)+(3-4t)>0.
\]
At $t=487/1000$ its value is the positive number
\[
 \frac{83847268664253747}{138198062500000000000},
\]
so a zero of $r_{v+5}$ has $t>487/1000$.
Equation~\eqref{eq:evenexception} then gives
\[
 E_{2v+5}<(1+t)^{-1}-tP_{4,4}(t)<0;
\]
the endpoint value is
$-46095940700109163/148700000000000000$, and the bound decreases thereafter.
This excludes the even pair for every $j<2v$.

Finally, $j<2v$ corresponds to $n\le5v+2$ in the odd case and to
$n\le6v$ in the even case.  Advancing by the relevant parity proves the two
stated lower bounds.
\end{proof}

The accompanying dependency-free checker verifies $11544$ two-block formulas,
$17390$ parity identities, and $23088$ predecessor-pair identities for all $74$
odd exponents through $v=151$.  An independent exact symbolic checker clears
denominators and uses rational real-root counts to certify the eight derivative
signs and two positivity assertions used above.  These finite audits support,
but are not logically substituted for, the uniform analytic proof.

\begin{theorem}[Three-block predecessor safe lattice]
\label{thm:threeblocks}
Let $0<t<1/2$, let $v\ge5$ be odd, and let $n>2v$.  If
\[
 q_{n-2}(t)=q_{n-v}(t)=0,
\]
then
\[
 n\ge7v+4\quad\hbox{when $n$ is odd},\qquad
 n\ge8v+2\quad\hbox{when $n$ is even}.
\]
Thus every first coefficient-one escape in
Proposition~\ref{prop:oneescape} obeys these stronger bounds.
\end{theorem}
\begin{proof}
Theorem~\ref{thm:safelattice} and~\eqref{eq:predecessorpairs} reduce the new
work to $j=2v+s$ with $0\le s<v$.  A third term in the geometric expansion
used for~\eqref{eq:firsttworblocks} gives
\begin{equation}\label{eq:thirdrblock}
 r_{2v+s}=P_{1,2v+s}+P_{3,v+s}+P_{5,s}.
\end{equation}
We use repeatedly the derivative identity
$P'_{k,m}=-(k+1)P_{k+1,m-1}$ and the paired-increment identity
\begin{equation}\label{eq:pairedincrement}
 P_{k,m+2}-P_{k,m}
 =(-1)^{m+1}{k+m+1\choose m+1}t^{m+1}
 \left(1-\frac{k+m+2}{m+2}t\right).
\end{equation}
Thus the relevant odd partial sums increase with their odd order and the
relevant even partial sums decrease with their even order.

We first isolate all possible zeros of~\eqref{eq:thirdrblock}.  If $s$ is even
and $s\ge2$, or if $s=0$ and $v\ge7$, Taylor's theorem and
\eqref{eq:pairedincrement} give
\[
 r_{2v+s}>(1+t)^{-2}+P_{3,7}(t)+(1+t)^{-6}>0.
\]
The last function is decreasing, and its value at $1/2$ is $4021/11664$.
For the remaining case $(v,s)=(5,0)$, use
$P_{3,5}(t)\ge P_{3,5}(1/2)=-9/16$ and
$P_{1,10}>(1+t)^{-2}\ge4/9$ to obtain
$r_{10}>127/144$.  If $s\ge13$ is odd, then $v\ge15$ and
\[
 r_{2v+s}>P_{1,43}(t)+(1+t)^{-4}+P_{5,13}(t)>0.
\]
Again the displayed function is decreasing; at $1/2$ it equals
\[
 \frac{108675557490085}{356241767399424}>0.
\]
Hence only $s\in\{1,3,5,7,9,11\}$ can remain.

For these six offsets let $v_s$ be the least admissible odd $v$, namely
$5,5,7,9,11,13$, and put
\[
 L_s=P_{1,2v_s+s}+(1+t)^{-4}+P_{5,s},\qquad
 U_s=(1+t)^{-2}+P_{3,v_s+s}+P_{5,s}.
\]
Odd partial sums increase and even partial sums decrease, so
$L_s<r_{2v+s}<U_s$.  Exact rational Sturm chains give the following brackets
for every zero of $r_{2v+s}$ in $(0,1/2)$:
\[
\begin{array}{c|cccccc}
s&1&3&5&7&9&11\\ \hline
\alpha_s&.3178&.3441&.3873&.4270&.4618&.4921\\
\beta_s &.322 &.345 &.388 &.428 &.462 &.493
\end{array}
\qquad \alpha_s<t<\beta_s.
\]
All decimals here are the displayed terminating rationals.  The certificate
checks that $L_s$ has no zero on $(0,\alpha_s]$ and is positive there, while
$U_s$ has no zero on $[\beta_s,1/2)$ and is negative there.  Their endpoint
margins are retained as exact fractions.

Consider first the odd event pair.  At a zero of $r_{2v+s}$, taking one more
partial-sum term in~\eqref{eq:thirdrblock} gives
\begin{align*}
 r_{2v+s+1}={}&(2v+s+2)t^{2v+s+1}
 -{v+s+4\choose3}t^{v+s+1}
 +{s+6\choose5}t^{s+1}.
\end{align*}
Dropping the negative middle term and using $t<\beta_s$, the remaining
$v$-dependent term is maximal at $v=v_s$: after replacing $v$ by $v+2$, its
ratio is at most
\[
 \frac{2v+s+6}{2v+s+2}\beta_s^4
 \le\frac{17}{13}\left(\frac{493}{1000}\right)^4<1.
\]
We therefore obtain
\[
 t r_{2v+s+1}<
 {s+6\choose5}\beta_s^{s+2}
 +(2v_s+s+2)\beta_s^{2v_s+s+2}<\frac34.
\]
The last six strict rational inequalities are recorded in the audit.  Thus
$E_{2v+s+2}=1-t r_{2v+s+1}>1/4$, excluding the odd pair.

For the even pair, the same geometric expansion gives, for every $m$,
\begin{equation}\label{eq:evenpartialblocks}
 E_m=P_{0,m}-t\sum_{1\le h\le m/v}P_{2h,m-hv-1}.
\end{equation}
At $m=3v+s$ this has the three terms $P_2,P_4,P_6$ in the sum.  If $s=1$,
the root bracket and partial-sum monotonicity imply
\[
 E_{3v+1}>(1+t)^{-1}
 -t\{P_{2,10}(t)+(1+t)^{-5}+1\}>0.
\]
The final function has no zero on $[.3178,.322]$ and is positive at both
endpoints.  For $s\in\{3,5,7,9,11\}$, use
$P_{0,3v+s}=(1+t^{3v+s+1})/(1+t)$, bound the positive tail by
$\beta_s^{3v_s+s+1}$, and use monotonicity in the other three partial sums:
\begin{align*}
 E_{3v+s}<&\ (1+t)^{-1}+\beta_s^{3v_s+s+1}\\
 &-t\{(1+t)^{-3}+P_{4,v_s+s-1}(t)+P_{6,s-1}(t)\}<0.
\end{align*}
For each $s$, the last upper bound has no zero on
$[\alpha_s,\beta_s]$ and is negative at both endpoints, again by exact
rational Sturm chains.  This excludes the even pair.

We have excluded both pairs for $j<3v$.  The largest corresponding odd event
index is $7v+2$, and the largest even event index is $8v$.  Advancing by the
relevant parity proves the theorem.
\end{proof}

The dependency-free checker verifies $30591$ three-block formulas,
$40887$ parity identities, and $61182$ predecessor-pair identities for all
$99$ odd exponents through $v=201$.  A separate symbolic checker supplies
$21$ exact root-free rational-interval sign certificates for the exceptional
offsets and companion bounds.  These audits are reproducible finite checks of
the explicit algebra in the uniform proof; the theorem is not inferred from
the cutoff $v=201$.

\begin{theorem}[Four-block predecessor safe lattice]
\label{thm:fourblocks}
Let $0<t<1/2$, let $v\ge5$ be odd, and let $n>2v$.  If
\[
 q_{n-2}(t)=q_{n-v}(t)=0,
\]
then
\[
 n\ge9v+4\quad\hbox{when $n$ is odd},\qquad
 n\ge10v+2\quad\hbox{when $n$ is even}.
\]
Thus every first coefficient-one escape in
Proposition~\ref{prop:oneescape} obeys these stronger bounds.
\end{theorem}
\begin{proof}
By Theorem~\ref{thm:threeblocks}, it remains to exclude the two pairs in
\eqref{eq:predecessorpairs} for $j=3v+s$, $0\le s<v$.  The next geometric
term gives
\begin{equation}\label{eq:fourthrblock}
 r_{3v+s}=P_{1,3v+s}+P_{3,2v+s}+P_{5,v+s}+P_{7,s}.
\end{equation}
We continue to use the paired-increment identity~\eqref{eq:pairedincrement}.

First isolate the zeros of~\eqref{eq:fourthrblock}.  Suppose that $s$ is
even.  For $s\in\{2,4,\ldots,16\}$, let $w_s=\max(5,s+1)$.  Odd partial
sums increase and even partial sums decrease toward their limits, so
\[
 r_{3v+s}>
 P_{1,3w_s+s}+(1+t)^{-4}+P_{5,w_s+s}+P_{7,s}>0.
\]
The last eight signs are exact Sturm certificates.  The same argument with
$w_0=9$ proves positivity for $s=0$ and $v\ge9$.  If $s\ge18$ is even, then
$P_{1,3v+s}>0$, while the $P_3$ and $P_7$ terms are positive even partial
sums.  Also $v+s\ge2s+1\ge37$, and
$P_{5,v+s}\ge P_{5,37}>0$; the last sign is again certified on $(0,1/2)$.
Thus the only even-offset cases not covered analytically are $(v,s)=(5,0)$
and $(7,0)$.

Now suppose that $s$ is odd.  For $s\ge19$, the least admissible exponent is
$v=s+2$, and monotonicity gives
\[
 r_{3v+s}>(1+t)^{-2}+P_{3,3s+4}+(1+t)^{-6}+P_{7,s}
 \ge (1+t)^{-2}+P_{3,61}+(1+t)^{-6}+P_{7,19}>0.
\]
The final function has no zero on $(0,1/2)$ and is positive at both
endpoints.  For $s\in\{1,3,\ldots,17\}$, put $v_1=7$ and
$v_s=\max(5,s+2)$ otherwise, and define
\begin{align*}
 L_s&=(1+t)^{-2}+P_{3,2v_s+s}+(1+t)^{-6}+P_{7,s},\\
 U_s&=P_{1,3v_s+s}+(1+t)^{-4}+P_{5,v_s+s}+P_{7,s}.
\end{align*}
Then $L_s<r_{3v+s}<U_s$.  Exact rational Sturm chains confine every zero to
the following intervals:
\[
\begin{array}{c|ccccc}
s&1&3&5&7&9\\ \hline
\alpha_s&.277&.288&.323&.359&.391\\
\beta_s &.279&.290&.325&.360&.392
\end{array}
\qquad
\begin{array}{c|cccc}
s&11&13&15&17\\ \hline
\alpha_s&.420&.446&.469&.490\\
\beta_s &.421&.447&.471&.492
\end{array}
\]
Here $L_s$ is positive and root-free on $(0,\alpha_s]$, while $U_s$ is
negative and root-free on $[\beta_s,1/2)$.  The only omitted odd case is
$(v,s)=(5,1)$.

The three finite cases $(5,0)$, $(7,0)$, and $(5,1)$ are disposed of without
approximating their roots.  In each case the two integer-polynomial pairs
$(r_{3v+s},E_{3v+s+2})$ and $(r_{3v+s},E_{4v+s})$ are primitive, their degrees
are preserved modulo $7$, and their monic gcds in $\mathbb F_7[t]$ are $1$.
The retained certificates give every Euclidean remainder.  Gauss's lemma
therefore excludes a common zero over $\C$, and hence over $(0,1/2)$.

It remains to separate the companions at the nine uniform root intervals.
At a zero of $r_{3v+s}$, the next partial-sum terms give
\begin{align*}
 r_{3v+s+1}={}&-(3v+s+2)t^{3v+s+1}
 +{2v+s+4\choose3}t^{2v+s+1}\\
 &-{v+s+6\choose5}t^{v+s+1}
 +{s+8\choose7}t^{s+1}.
\end{align*}
After multiplication by $t$, discard the two negative terms.  The ratio of
the remaining $v$-dependent term after replacing $v$ by $v+2$ is at most
\[
 \frac{{2v+s+8\choose3}}{{2v+s+4\choose3}}\,\beta_s^4<1;
\]
it is enough to check this at $v=v_s$.  The nine exact rational margins give
\[
 t r_{3v+s+1}<
 {2v_s+s+4\choose3}\beta_s^{2v_s+s+2}
 +{s+8\choose7}\beta_s^{s+2}<1.
\]
Consequently $E_{3v+s+2}=1-t r_{3v+s+1}>0$, excluding the odd event pair.

For the even event pair, substitute $m=4v+s$ in
\eqref{eq:evenpartialblocks}.  When $s=1$, partial-sum monotonicity and the
root bracket give
\begin{align*}
 E_{4v+1}>&\ \frac{1-\beta_1^{4v_1+2}}{1+t}\\
 &-t\{(1+t)^{-3}+P_{4,2v_1}(t)+(1+t)^{-7}+1\}>0.
\end{align*}
For $s\in\{3,5,\ldots,17\}$, the complementary bounds give
\[
 E_{4v+s}< (1+t)^{-1}
 -t\{P_{2,3v_s+s-1}(t)+(1+t)^{-5}
       +P_{6,v_s+s-1}(t)+P_{8,s-1}(t)\}<0.
\]
Each displayed companion bound is root-free with the stated sign on
$[\alpha_s,\beta_s]$, by exact Sturm chains.  This excludes the even event
pair.

We have excluded both pairs for $j<4v$.  The largest corresponding odd event
index is $9v+2$, and the largest even event index is $10v$.  Advancing by the
relevant parity proves the theorem.
\end{proof}

The dependency-free checker verifies $63488$ four-block formulas, $79484$
parity identities, and $126976$ predecessor-pair identities for all $124$ odd
exponents through $v=251$.  It also retains full mod-$7$ Euclidean
certificates for the six finite exceptional companion pairs.  An independent
SymPy checker supplies $38$ root-free rational-interval sign certificates and
recomputes all six modular gcds.  Both deterministic JSON outputs use
canonical LF bytes, making their recorded hashes platform-independent.  The
finite cutoff audits the algebra; the theorem itself is uniform.

\begin{theorem}[Five-block predecessor safe lattice]\label{thm:fiveblocks}
Let $0<t<1/2$, let $v\ge5$ be odd, and let $n>2v$.  If
\[
 q_{n-2}(t)=q_{n-v}(t)=0,
\]
then
\[
 n\ge11v+4\quad\hbox{when $n$ is odd},\qquad
 n\ge12v+2\quad\hbox{when $n$ is even}.
\]
Thus every first coefficient-one escape in Proposition~\ref{prop:oneescape}
obeys these stronger bounds.
\end{theorem}
\begin{proof}
By Theorem~\ref{thm:fourblocks}, it remains to exclude the two pairs in
\eqref{eq:predecessorpairs} for $j=4v+s$, $0\le s<v$.  The next geometric
term gives
\begin{equation}\label{eq:fiveblock}
 r_{4v+s}=P_{1,4v+s}+P_{3,3v+s}+P_{5,2v+s}
              +P_{7,v+s}+P_{9,s}.
\end{equation}

We first isolate its zeros.  Suppose that $s$ is even.  For
$s\in\{0,2,4,6\}$ use respectively
\[
 (s,v_0)=(0,17),(2,7),(4,5),(6,7)
\]
and put
\[
 M_{s,v_0}=(1+t)^{-2}+P_{3,3v_0+s}+(1+t)^{-6}
                         +P_{7,v_0+s}+P_{9,s}.
\]
Even partial sums lie above their limits and odd partial sums increase with
their odd order.  Hence $r_{4v+s}>M_{s,v_0}$ whenever $v\ge v_0$.
Exact rational Sturm chains show that these four functions are positive and
root-free on $(0,1/2)$.  This leaves only
\begin{equation}\label{eq:fiveblockevenexceptions}
 (v,s)=(5,0),(5,2),(7,0),(9,0),(11,0),(13,0),(15,0).
\end{equation}

For every even $s\ge8$, the same monotonicity gives
\[
 r_{4v+s}>M_s:=(1+t)^{-2}+P_{3,4s+3}+(1+t)^{-6}
                         +P_{7,2s+1}+P_{9,s}.
\]
Write $A_9={s+10\choose9}$, $A_7={2s+9\choose7}$, and
$A_3={4s+7\choose3}$.  Alternating-tail estimates give
\begin{align*}
 P_{9,s}-(1+t)^{-10}
 &\ge A_9t^{s+1}\left(1-\frac{s+11}{s+2}t\right),\\
 (1+t)^{-8}-P_{7,2s+1}&<A_7t^{2s+2},\\
 (1+t)^{-4}-P_{3,4s+3}&<A_3t^{4s+4}.
\end{align*}
It is enough that $R_7(s)+R_3(s)<1$, where
\[
 R_7(s)=\frac{A_7}{A_9}\frac{2(s+2)}{s-7}\,2^{-s-1},\qquad
 R_3(s)=\frac{A_3}{A_9}\frac{2(s+2)}{s-7}\,2^{-3s-3}.
\]
At $s=8$ their sum is $126012629939/326283296768<1$.  Both summands
decrease under $s\mapsto s+2$: their successive ratios are
\begin{align*}
 &\frac{(s-7)(s+4)(s+5)(s+6)(2s+11)(2s+13)}
 {4(s-5)(s+2)(s+11)(s+12)(2s+3)(2s+5)},\\
 &\frac{(s-7)(s+3)(s+4)(2s+7)(4s+13)(4s+15)}
 {64(s-5)(s+11)(s+12)(2s+3)(4s+5)(4s+7)}.
\end{align*}
After $s=x+8$, the numerators of one minus these ratios have only positive
coefficients.  Thus $M_s>0$, excluding every remaining even offset.

Now let $s$ be odd.  Put $v_1=13$, and for
$s\in\{3,5,\ldots,23\}$ put $v_s=\max(5,s+2)$.  Define
\begin{align*}
 L_s={}&P_{1,4v_s+s}+(1+t)^{-4}+P_{5,2v_s+s}
                         +(1+t)^{-8}+P_{9,s},\\
 U_s={}&(1+t)^{-2}+P_{3,3v_s+s}+(1+t)^{-6}
                         +P_{7,v_s+s}+P_{9,s}.
\end{align*}
For $v\ge v_s$, the paired-increment identity gives
$L_s<r_{4v+s}<U_s$.  For odd $s\ge25$, the lower bound increases with $s$,
and its first instance
\[
 P_{1,133}+(1+t)^{-4}+P_{5,79}+(1+t)^{-8}+P_{9,25}
\]
is positive and root-free on $(0,1/2)$ by an exact Sturm chain.  Hence only
the twelve offsets $s=1,3,\ldots,23$ remain.  Exact rational Sturm chains
confine their uniform zeros to
\[
\begin{array}{c|cccccc}
s&1&3&5&7&9&11\\ \hline
\alpha_s&.248&.249&.279&.310&.340&.367\\
\beta_s &.250&.251&.280&.312&.341&.368
\end{array}
\qquad
\begin{array}{c|cccccc}
s&13&15&17&19&21&23\\ \hline
\alpha_s&.392&.415&.435&.454&.472&.488\\
\beta_s &.393&.416&.437&.456&.474&.490.
\end{array}
\]
Here $L_s$ is positive and root-free on $(0,\alpha_s]$, while $U_s$ is
negative and root-free on $[\beta_s,1/2)$.  At $s=1$ this uses $v\ge13$;
the omitted cases are
\begin{equation}\label{eq:fiveblockoddexceptions}
 (v,s)=(5,1),(7,1),(9,1),(11,1).
\end{equation}

All eleven cases in \eqref{eq:fiveblockevenexceptions} and
\eqref{eq:fiveblockoddexceptions} are disposed of without approximating a
root.  For each case both integer-polynomial pairs
\[
 (r_{4v+s},E_{4v+s+2}),\qquad (r_{4v+s},E_{5v+s})
\]
are primitive and have degree-preserving reduction with monic gcd $1$ over a
  finite field.  Modulus $7$ works except for
  $(v,s,\mathrm{companion})=(5,0,\mathrm{odd}),
  (5,0,\mathrm{even}),(5,2,\mathrm{even})$, for which modulus $11$ works.
  The retained certificates contain every Euclidean remainder.
Gauss's lemma excludes a common zero over $\C$.

It remains to separate the companions on the twelve uniform intervals.  At a
zero of $r_{4v+s}$, the five next partial-sum terms give
\begin{align*}
 r_{4v+s+1}={}&(4v+s+2)t^{4v+s+1}
 -{3v+s+4\choose3}t^{3v+s+1}\\
 &+{2v+s+6\choose5}t^{2v+s+1}
 -{v+s+8\choose7}t^{v+s+1}
 +{s+10\choose9}t^{s+1}.
\end{align*}
After multiplication by $t$, discard the two negative terms.  The two
$v$-dependent positive terms decrease under $v\mapsto v+2$ on the root
interval because
\[
 \frac{4(v_s+2)+s+2}{4v_s+s+2}\,\beta_s^8<1,\qquad
 \frac{{2(v_s+2)+s+6\choose5}}{{2v_s+s+6\choose5}}\,\beta_s^4<1.
\]
The twelve exact margins give
\[
 (4v_s+s+2)\beta_s^{4v_s+s+2}
 +{2v_s+s+6\choose5}\beta_s^{2v_s+s+2}
 +{s+10\choose9}\beta_s^{s+2}<1.
\]
Consequently $E_{4v+s+2}=1-tr_{4v+s+1}>0$.

For the even event pair, substitute $m=5v+s$ in
\eqref{eq:evenpartialblocks}.  At $s=1$, partial-sum monotonicity gives
\[
 E_{5v+1}>(1+t)^{-1}
 -t\{P_{2,4v_1}+(1+t)^{-5}+P_{6,2v_1}+(1+t)^{-9}+1\}>0.
\]
For $s\in\{3,5,\ldots,23\}$, complementary bounds give
\begin{align*}
 E_{5v+s}<{}&\frac{1+t^{5v_s+s+1}}{1+t}\\
 &-t\{(1+t)^{-3}+P_{4,3v_s+s-1}+(1+t)^{-7}
             +P_{8,v_s+s-1}+P_{10,s-1}\}<0.
\end{align*}
Each displayed bound is root-free with the stated sign on
$[\alpha_s,\beta_s]$, by exact rational Sturm chains.

We have excluded both pairs for $j<5v$.  The largest corresponding odd event
index is $11v+2$, and the largest even event index is $12v$.  Advancing by
the relevant parity proves the theorem.
\end{proof}

The dependency-free checker verifies $79360$ five-block formulas, $95356$
parity identities, and $158720$ predecessor-pair identities for all $124$
odd exponents through $v=251$, together with full modular Euclidean
certificates for the $22$ finite exceptional pairs.  An independent SymPy
checker supplies $41$ root-free rational-interval sign certificates and
recomputes all $22$ modular gcds.  A separate exact reconnaissance through
$v=47$ records the discovery pattern but is not used to infer the theorem.
All deterministic outputs use canonical LF bytes; the finite audit windows
check the algebra, while the proof above is uniform.

\begin{theorem}[Six-block predecessor safe lattice]\label{thm:sixblocks}
Let $0<t<1/2$, let $v\ge5$ be odd, and let $n>2v$.  If
\[
 q_{n-2}(t)=q_{n-v}(t)=0,
\]
then
\[
 n\ge13v+4\quad\hbox{when $n$ is odd},\qquad
 n\ge14v+2\quad\hbox{when $n$ is even}.
\]
Thus every first coefficient-one escape in Proposition~\ref{prop:oneescape}
obeys these stronger bounds.
\end{theorem}
\begin{proof}
By Theorem~\ref{thm:fiveblocks}, it remains to exclude
\eqref{eq:predecessorpairs} for $j=5v+s$, $0\le s<v$.  The sixth geometric
term gives
\begin{equation}\label{eq:sixblock}
 r_{5v+s}=P_{1,5v+s}+P_{3,4v+s}+P_{5,3v+s}
             +P_{7,2v+s}+P_{9,v+s}+P_{11,s}.
\end{equation}

We first exclude even $s$.  For
\[
 (s,v_0)=(0,23),(2,15),(4,5),(6,7),(8,9)
\]
put
\[
 M_{s,v_0}=P_{1,5v_0+s}+(1+t)^{-4}+P_{5,3v_0+s}
 +(1+t)^{-8}+P_{9,v_0+s}+P_{11,s}.
\]
Odd partial sums increase with their odd order and even partial sums lie above
their limits.  Hence $r_{5v+s}>M_{s,v_0}$ for $v\ge v_0$.  Exact rational
Horner interval enclosures certify $M_{s,v_0}>0$ on $(0,1/2)$ in all five
cases.  The omitted pairs are
\begin{equation}\label{eq:sixblockevenexceptions}
 \{(v,0):5\le v\le21\}\cup\{(v,2):5\le v\le13\},
 \qquad v\ \hbox{odd}.
\end{equation}

For every even $s\ge10$, take $v_0=s+1$.  The same lower bound applies.
Write
\[
 A_{11}={s+12\choose11},\quad A_9={2s+11\choose9},\quad
 A_5={4s+9\choose5},\quad A_1=6s+7.
\]
The alternating remainder bounds compare the positive remainder of
$P_{11,s}$ with the three negative remainders:
\begin{align*}
 P_{11,s}-(1+t)^{-12}
 &\ge A_{11}t^{s+1}\left(1-\frac{s+13}{s+2}t\right),\\
 (1+t)^{-10}-P_{9,2s+1}&<A_9t^{2s+2},\\
 (1+t)^{-6}-P_{5,4s+3}&<A_5t^{4s+4},\\
 (1+t)^{-2}-P_{1,6s+5}&<A_1t^{6s+6}.
\end{align*}
It is enough that
\[
 \frac{2(s+2)}{s-9}\frac1{A_{11}}
 \left(\frac{A_9}{2^{s+1}}+\frac{A_5}{2^{3s+3}}
                         +\frac{A_1}{2^{5s+5}}\right)<1.
\]
At $s=10$ the left side is
$27281522928263929231/81461110059877531648<1$.  Each of its three
summands decreases under $s\mapsto s+2$; after $s=x+10$, the exact
numerators of one minus the successive ratios have only positive
coefficients.  This excludes every even offset outside
\eqref{eq:sixblockevenexceptions}.

Now let $s$ be odd.  Put $v_1=19$, $v_3=11$, and
$v_s=s+2$ for $s\ge5$.  Define
\begin{align*}
 L_s={}&(1+t)^{-2}+P_{3,4v_s+s}+(1+t)^{-6}
                  +P_{7,2v_s+s}+(1+t)^{-10}+P_{11,s},\\
 U_s={}&P_{1,5v_s+s}+(1+t)^{-4}+P_{5,3v_s+s}
                  +(1+t)^{-8}+P_{9,v_s+s}+P_{11,s}.
\end{align*}
Paired increments give $L_s<r_{5v+s}<U_s$ for $v\ge v_s$.
For odd $s\ge33$, the three negative Taylor remainders in $L_s$ are bounded
at $t=1/2$ by
\[
 {s+12\choose11}2^{-s-1}+{3s+12\choose7}2^{-3s-5}
                       +{5s+12\choose3}2^{-5s-9}.
\]
At $s=33$ this is
\[
 \frac{1768483226579352197428635688881857408474171792538535}
 {2993155353253689176481146537402947624255349848014848}
 <\frac{421876}{531441}
 =\sum_{k=0}^{5}\left(\frac23\right)^{2k+2}.
\]
All three errors decrease under $s\mapsto s+2$, again by positive-coefficient
ratio certificates.  Hence $L_s>0$ for every odd $s\ge33$.

The remaining offsets $s=1,3,\ldots,31$ are confined to the following
rational intervals:
\[
\begin{array}{c|rrrrrrrr}
s&1&3&5&7&9&11&13&15\\ \hline
\alpha_s&.226&.220&.246&.274&.301&.326&.350&.371\\
\beta_s &.228&.222&.248&.276&.303&.329&.352&.374
\end{array}
\]
\[
\begin{array}{c|rrrrrrrr}
s&17&19&21&23&25&27&29&31\\ \hline
\alpha_s&.391&.410&.427&.443&.458&.472&.486&.499\\
\beta_s &.394&.413&.430&.446&.462&.476&.489&.4999
\end{array}
\]
Exact rational interval Horner certificates show $L_s>0$ on
$(0,\alpha_s]$ and $U_s<0$ on $[\beta_s,1/2)$.  The cases omitted by the
choices of $v_1$ and $v_3$ are
\begin{equation}\label{eq:sixblockoddexceptions}
 \{(v,1):5\le v\le17\}\cup\{(v,3):5\le v\le9\},
 \qquad v\ \hbox{odd}.
\end{equation}

Every one of the $24$ cases in
\eqref{eq:sixblockevenexceptions}--\eqref{eq:sixblockoddexceptions} is
disposed of algebraically.  For both pairs
\[
 (r_{5v+s},E_{5v+s+2}),\qquad(r_{5v+s},E_{6v+s}),
\]
the integer polynomials are primitive and have degree-preserving reduction
with monic gcd $1$ over one of $\mathbb F_2,\mathbb F_3,\mathbb F_5,
\mathbb F_7$.  The retained $48$ certificates contain every Euclidean
remainder; Gauss's lemma excludes a common complex zero.

It remains to separate the companions on the $16$ uniform intervals.  At a
zero of $r_{5v+s}$, the six next terms alternate.  Discarding the three
negative terms gives
\begin{align*}
 t r_{5v+s+1}<&{4v+s+4\choose3}t^{4v+s+2}
 +{2v+s+8\choose7}t^{2v+s+2}
 +{s+12\choose11}t^{s+2}.
\end{align*}
The two $v$-dependent terms decrease under $v\mapsto v+2$ on each root
interval, and their exact endpoint sum is strictly below $1$ at $v=v_s$.
Thus $E_{5v+s+2}=1-tr_{5v+s+1}>0$.

Finally, substitute $m=6v+s$ in~\eqref{eq:evenpartialblocks}.  For $s=1$,
partial-sum monotonicity gives the positive lower bound
\begin{align*}
 E_{6v+1}>{}&\frac{1-t^{6v_1+2}}{1+t}
 -t\{(1+t)^{-3}+P_{4,4v_1}+(1+t)^{-7}\\
 &\hspace{43mm}+P_{8,2v_1}+(1+t)^{-11}+1\}>0.
\end{align*}
For $s=3,5,\ldots,31$, complementary bounds give
\begin{align*}
 E_{6v+s}<&\frac1{1+t}-t\{P_{2,5v_s+s-1}+(1+t)^{-5}
 +P_{6,3v_s+s-1}\\
 &\hspace{34mm}+(1+t)^{-9}+P_{10,v_s+s-1}+P_{12,s-1}\}<0.
\end{align*}
All $16$ signs hold throughout $[\alpha_s,\beta_s]$ by the retained exact
interval certificates.  Hence both predecessor pairs are excluded for
$j<6v$.  The largest odd and even event indices are $13v+2$ and $14v$;
advancing by parity proves the theorem.
\end{proof}

The dependency-free checker verifies $95232$ six-block formulas, $111228$
parity identities, and $190464$ predecessor-pair identities for all $124$ odd
exponents through $v=251$, together with full modular Euclidean certificates
for the $48$ finite pairs.  An independent SymPy checker supplies $53$ exact
rational interval-Horner sign certificates and recomputes all $48$ modular
gcds.  Both theorem-level JSON outputs replay byte-for-byte.  Exact finite
reconnaissance through $v=23$ records $89$ roots and constant companion gcds;
it is discovery evidence only and is not used in the uniform proof.

\begin{remark}[An exact open bottleneck]\label{rem:coprimebottleneck}
Regard the $q_j$ as polynomials in an indeterminate $t$.  Exact modular gcd
certificates accompanying this paper prove
\[
 \gcd_{\Z[t]}(q_{n-2}(t),q_{n-v}(t))=1
\]
for all $876456$ pairs with odd $5\le v\le151$ and $2v<n\le12000$, with no
unresolved index.  This is finite evidence, not an infinite
proof.  Theorem~\ref{thm:sixblocks} independently removes every real event
below its parity-dependent linear thresholds.  Proving the displayed
coprimality for every odd $v\ge5$ and $n>2v$
would eliminate the only branch identified by
Proposition~\ref{prop:oneescape}.  Together with
Lemma~\ref{lem:pringsheim}, it would settle the entire normalized
smaller-gap-$2$ trinomial family, without restricting the algebraic degree of
$a$.
\end{remark}

We now use the partial sums above to give a finite exact representation of all
coefficients needed in Lemma~\ref{lem:forcedseries}.
Expanding~\eqref{eq:forcedseries} first in powers of $x^v$ gives
\begin{equation}\label{eq:blockexpansion}
 Q_{a,v}(x)=\frac1{(1-x^2)(1+a x^2)}+
 \sum_{k\ge1}\frac{(-1)^{k-1}a x^{kv+2}}
 {(1-x^2)(1+a x^2)^{k+1}}.
\end{equation}
Since $v$ is odd, equal parities combine.  For $m\ge0$ this yields
\begin{align}
 q_{Kv+2+2m}
 &=a\sum_{s=0}^{(K-1)/2}P_{K-2s,m+sv}(a), &&K\text{ odd},
 \label{eq:oddblock}\\
 q_{Kv+2+2m}
 &=P_{0,(K/2)v+1+m}(a)
   -a\sum_{s=0}^{K/2-1}P_{K-2s,m+sv}(a), &&K\text{ even}.
 \label{eq:evenblock}
\end{align}
These identities are formal and exact.
Keeping the top partial sum fixed and sending odd $v$ to infinity gives
\begin{equation}\label{eq:blocklimit}
 L_{K,m}(a)=
 \begin{cases}
 aP_{K,m}(a)+a\displaystyle\sum_{\substack{1\le j<K\\j\ \mathrm{odd}}}
 (1+a)^{-j-1},&K\text{ odd},\\[2mm]
 (1+a)^{-1}-aP_{K,m}(a)
 -a\displaystyle\sum_{\substack{2\le j<K\\j\ \mathrm{even}}}
 (1+a)^{-j-1},&K\text{ even}.
 \end{cases}
\end{equation}

For the uniform estimates, if $0<a\le U<1/2$ and
\[
 U\frac{k+m+2}{m+2}<1,
\]
then the alternating terms after $P_{k,m}$ decrease in magnitude, and
\begin{equation}\label{eq:tailbound}
 \left|P_{k,m}(a)-(1+a)^{-k-1}\right|
 \le {k+m+1\choose m+1}U^{m+1}.
\end{equation}
Thus every block has an exact $v\to\infty$ limit, and all terms omitted from
that limit have explicit rational bounds.  Increasing odd $v$ only decreases
those bounds once~\eqref{eq:tailbound} applies.

\begin{theorem}[Global cubic exclusion]\label{thm:globalcubic}
Let $C=A_aB$ be a normalized nonnegative factorization of a Newman polynomial,
where
\[
 A_a(x)=1+a x^2+x^v,\qquad v>2,\qquad a\notin\Z.
\]
Then $[\Q(a):\Q]\ne3$.  Consequently, in the smaller-gap-$2$ case every
genuinely fractional coefficient has algebraic degree at least $4$, without
any restriction on the product degree.
\end{theorem}
\begin{proof}
Theorem~\ref{thm:divisible} removes even $v$ and the boundary cases; after
reversal the remaining cases have odd $v\ge5$.  Lemma~\ref{lem:cubiccatalog}
reduces a cubic coefficient to $126$ exact minimal-polynomial records.

For $v=5,7,\ldots,21$, the verifier works directly in
$\Z[a]/(m_a)$.  Multiplication by $a$ and the recurrence
\[
 q_n=1-aq_{n-2}-q_{n-v}
\]
produce a triple of integers for every coefficient.  Rational bisection of
the unique target root in $(0,1/2)$ separates every sign.  In all $1134$
catalog--case pairs, the only zero before the first negative coefficient is
$q_v$ (besides the initial odd zero ray).  The verifier also certifies the
forcing term $a q_{n-2}+q_{n-v}$ strictly positive at every step, so each
preceding coefficient is strictly less than $1$.

For every odd $v\ge23$, formulas~\eqref{eq:oddblock} and
\eqref{eq:evenblock} give a uniform certificate.  Each catalog entry records
one obstruction index
\[
\begin{gathered}
 N=Kv+r,\qquad K\le27,\\
 (K,r)\in\{(K,4):K\text{ odd}\}\cup
 \{(K,6):K\text{ even}\}\cup
 \{(K,8):K\text{ odd}\}.
\end{gathered}
\]
For a fixed parity, the ranges $0\le m\le v-1$ in consecutive values of $K$
partition all coefficients after the initial rays.  Proving a complete block
positive for every $m\ge0$ is therefore stronger than the required statement.
The base block $K=0$ is elementary:
\[
 q_{2+2m}=P_{0,m+1}(a)=\frac{1-(-a)^{m+2}}{1+a}\in(0,1)
 \qquad(0\le m\le v-1,\ 0<a<1).
\]
For every block $K\ge1$, using an outward dyadic interval containing the
target root, the verifier checks every internal block offset $0\le m\le40$ by
rational interval arithmetic.  It then uses~\eqref{eq:tailbound} to prove all
later offsets positive.  All complete blocks before $K$ are positive, as is
the prefix of block $K$ before $r$; the only preceding noninitial zero is
$q_v=0$.
Finally, the same tail estimate proves $q_{Kv+r}<0$ uniformly from $v=23$
onward.  The weakest retained rational upper margin is less than
$-2.29\cdot10^{-4}$, while the weakest pre-obstruction lower margin is greater
than $5.19\cdot10^{-4}$.  These decimals merely summarize exact fractions.

An independent identity checker compares $29193$ coefficients obtained by
direct formal division with the block expansion at $v=23$.  A second run at
\mbox{$v=99$} compares $123053$ identities.  A separate rational-specialization
checker, independent of the cubic quotient arithmetic, compares $9009$ more
identities.  A stress checker iterates the algebraic recurrence directly for
all $4914$ pairs with odd $23\le v\le99$, without using the block-limit
inequalities.  These audits are not needed for the infinite conclusion.
Lemma~\ref{lem:forcedseries} now supplies the contradiction for every catalog
entry and every odd $v\ge5$.  Combining with
Theorem~\ref{thm:noquadratic} proves the asserted degree-at-least-four
conclusion.
\end{proof}

\section{Global quartic exclusion for the smaller gap two}
\label{sec:globalquartic}

The same forced-series mechanism extends one algebraic degree further, but the
finite arithmetic layer is substantially larger.  The catalog below is
exhaustive under the conjugate bounds, and every retained sign is certified by
integer or rational arithmetic.

\begin{lemma}[Uniform quartic catalog]\label{lem:quarticcatalog}
Suppose that $a\in(0,1/2)$ is a quartic algebraic integer permitted by the
norm and conjugate restrictions of the smaller-gap-$2$ case.  Then its monic
minimal polynomial
\[
 m_a(T)=T^4+s_3T^3+s_2T^2+s_1T+s_0
\]
belongs to an explicit catalog of $6{,}977$ irreducible integer polynomials.
Five catalog polynomials have two roots in $(0,1/2)$ and all others have one,
so the catalog contains $6{,}982$ target roots.
\end{lemma}
\begin{proof}
Every conjugate has modulus strictly less than $R=143/50$, every real
conjugate is greater than $-2$, and the distinguished root satisfies
$0<a<1/2$.  Vieta's formulas therefore give the integer bounds
\[
 |s_3|\le9,\qquad |s_2|\le28,\qquad |s_1|\le35,
 \qquad |s_0|\le11.
\]
For example, $|s_3|<1/2+3R$, and the other inequalities follow by separating
the elementary symmetric terms that contain $a$ from those that do not.  The
resulting box contains $1{,}691{,}646$ tuples with $s_0\ne0$.

The verifier applies the strict Schur recursion to
$50^4m_a((143/50)T)$, leaving $63{,}546$ polynomials whose roots all lie in
$|z|<R$.  Exact integer linear- and quadratic-factor tests leave $57{,}238$
irreducibles.  Rational Sturm chains retain $10{,}671$ polynomials having a
root in $(0,1/2)$ and then remove every polynomial with a real root in
$(-3,-2]$; the disk bound makes this the entire forbidden real range.  The
remaining counts are exactly those stated.  The same Sturm chains isolate
every target root in a rational interval of width at most $2^{-36}$.
\end{proof}

\begin{theorem}[Global quartic exclusion]\label{thm:globalquartic}
Let $C=A_aB$ be a normalized nonnegative factorization of a Newman polynomial,
where
\[
 A_a(x)=1+a x^2+x^v,\qquad v>2,\qquad a\notin\Z.
\]
Then $[\Q(a):\Q]\ne4$.  Consequently, in the smaller-gap-$2$ case every
genuinely fractional coefficient has algebraic degree at least $5$, without
any restriction on the product degree.
\end{theorem}
\begin{proof}
As in Theorem~\ref{thm:globalcubic}, it is enough to consider odd $v\ge5$ and
the target roots in Lemma~\ref{lem:quarticcatalog}.  For each of the $6{,}982$
roots and every $v=5,7,\ldots,21$, the verifier iterates
\[
 q_n=1-aq_{n-2}-q_{n-v}
\]
exactly in $\Z[a]/(m_a)$.  Rational interval evaluation separates every
algebraic sign.  Across all $62{,}838$ root--exponent pairs, the only
noninitial zero before the first negative coefficient is $q_v$, and every
forcing term preceding that negative coefficient is strictly positive.

For odd $v\ge23$, formulas~\eqref{eq:oddblock}--\eqref{eq:tailbound} give a
uniform certificate.  The expanded obstruction search uses every
parity-compatible even offset through $100$: offsets divisible by $4$ for odd
blocks and offsets congruent to $2$ modulo $4$ for even blocks.  Each target
root has a certified obstruction
\[
 N=Kv+r,\qquad K\le110,\qquad r\le100,
\]
whose exact limiting upper bound plus its finite-$v$ error is negative.  This
expanded offset set is essential: restricting to offsets $4$, $6$, and $8$
can skip an earlier negative coefficient at offset $10$.

For every selected obstruction, a separate exact verifier proves all complete
preceding blocks positive and also the prefix of the obstruction block before
$r$.  It evaluates internal offsets $0\le m\le48$ by rational interval
arithmetic and uses the full-block limit with a first-omitted alternating-tail
bound for every later offset.  The base block is positive by the formula in
the proof of Theorem~\ref{thm:globalcubic}.  At every intervening index at
least one of the two predecessor coefficients is therefore positive; hence
$a q_{n-2}+q_{n-v}>0$ and $q_n<1$.  Thus the hypotheses of
Lemma~\ref{lem:forcedseries} hold up to $N$, where the forced coefficient is
negative.  This contradicts the existence of the normalized nonnegative
cofactor and excludes every quartic target root.
\end{proof}

The quartic proof objects retain the ordered minimal-polynomial catalog, both
target intervals for the five double-target polynomials, every finite
algebraic sign record, one exact uniform negative-tail certificate per target
root, and the complete positive-prefix audit.  Floating point is used only to
order candidate obstruction checks; every accepted inequality is recomputed
with exact rational endpoints.

For a final independent audit, a second implementation replaces the catalog
generator's Schur recursion by the Schur--Cohn Hermitian determinant criterion,
factors every survivor independently with FLINT, and encloses all complex roots
with Arb.  It reproduces the ordered catalog and independently validates all
$6{,}982$ retained target intervals.  At $4096$-bit ball precision, a
separate direct scalar recurrence checks all $62{,}838$ finite root--exponent
pairs, and an independent
parity-block implementation checks every retained uniform tail and positive
prefix certificate.  In total these two sign audits certify $12{,}369{,}079$
finite coefficient inequalities, $6{,}982$ uniform negative-tail
certificates, and $4{,}979{,}348$ uniform prefix inequalities, with no
indeterminate or failed sign.  These audits are redundant: the exact
rational proof objects remain the proof data used above.

\section{Literature boundary and next problems}

The global degree-exclusion theorems have to be separated from several
adjacent results.
Ghidelli proves fairness for $1+a x^2+x^5$ with an arbitrary nonnegative
cofactor and no product-degree cap~\cite{Ghidelli}.  Consequently the $v=5$
parts of Theorems~\ref{thm:globalcubic} and~\ref{thm:globalquartic} are
re-verifications in algebraic degrees $3$ and $4$, not new mathematical cases.
Ghidelli's proof uses linear recurrences, nonvanishing resultants, and root
asymptotics for that fixed exponent.  The present parity-block expansion uses
the arithmetic discreteness of the cubic and quartic catalogs to cover every
odd exponent uniformly.  Hare's algorithm instead scans candidate-factor
coefficient patterns of factor degree at most $15$ and eliminates all but
$975$ among $7{,}141{,}686$ patterns~\cite{Hare}.  That factor-degree
classification and
the present unbounded-exponent but fixed algebraic-degree theorems constrain
different parameters, so neither computation should be described as subsuming
the other.

There is also a direct integer-divisor literature.  Drungilas, Jankauskas, and
\v{S}iurys give a finite-state algorithm deciding restricted-coefficient
multiples for monic integer polynomials without unit-circle zeros and apply it
to low-degree Borwein polynomials~\cite{DJS}.  Idris and Sac-\'Ep\'ee study
which integer polynomials divide Newman polynomials, including explicit
nondivisibility and bounded-multiple computations~\cite{IS}.  Their subsequent
work gives additional existence and nonexistence results for Newman and
Littlewood multiples of certain integer polynomials~\cite{IS2}.  These methods
become relevant here after the conjugate product $H_a$ converts a nonnegative
real-factor question into an integer Newman-divisor question.  The global cubic
and quartic proofs instead return to the distinguished real embedding and
certify the forced cofactor recurrence directly.  The literature sources
checked did not exhibit the parity-block certificate or an all-odd-$v$ quartic
exclusion, but that absence is a literature-search inference rather than a
priority proof.

Up to reversal take $1\le u\le v/2$.  Theorem~\ref{thm:divisible} resolves
precisely the classes with $u\mid v$.  For every remaining class,
Theorem~\ref{thm:half} removes $[1/2,1)$ and Theorem~\ref{thm:integral}
removes rational points.  The norm sieve converts every fixed algebraic degree
into a finite search.

Three directions remain concrete.  First, Lemma~\ref{lem:pringsheim} already
gives an eventual negative coefficient for every real $a>0$; to turn that into
full fairness for the smaller gap two, push the six-block safe lattice beyond
$j=6v$ or rule out every nonstructural zero pair that could let a finite forced
cofactor terminate with leading coefficient $1$ before the first negative
term.
Second, determine whether the norm sieve and parity-block method remain
computationally and structurally effective in algebraic degree $5$.  Third,
organize propagation for smaller gaps
$u>2$ by residue classes or the continued fraction of $v/u$, and compare
Theorem~\ref{thm:norm} systematically with general
results on prescribed Newman divisors, especially the algorithms of Idris and
Sac-\'Ep\'ee~\cite{IS}.  The current paper proves neither the full conjecture
nor a counterexample.

\section*{AI research provenance}

This manuscript results from a state-aware autonomous AI Researcher--Critic
workflow.  During each research iteration, the AI Researcher explored proofs,
derived symbolic identities, searched for counterexamples, ran exact
computations, implemented verification programs, surveyed relevant literature,
and drafted the manuscript.  In separate review iterations, the AI Critic
checked claims, replayed computations, searched for gaps and edge cases, and
returned structured feedback for the next Researcher iteration.

The human role was limited to initiating the project, defining the overall
objective, configuring and supervising the workflow, selecting retained
outputs, and authorizing release.  The proof arguments, computational searches,
verification code, critical assessments, and prose were generated primarily by
the AI workflow and were not independently derived by the human orchestrator.
The labels ``Researcher'' and ``Critic'' describe workflow roles rather than
additional authors.

The accompanying repository provides the manuscript source, exact
computational certificates, verification programs, and integrity manifests so
that the workflow's mathematical outputs can be inspected and reproduced.
The complete public workflow record, including the stable automation prompt,
schedule, version chronology, and retained timing ledger, is provided in
\texttt{docs/autonomous-research-workflow.md}.

\subsection*{Prompted workflow}

The automation used a two-role, state-aware cycle.  The Researcher prompt
required at least 300 active minutes of proof work, exact computation,
counterexample search, and literature checking, with exactly one incremented
manuscript version created only after material progress.  The Critic prompt
required at least 60 active minutes of separate review, including computation
replay, gap and edge-case searches, and severity-ranked feedback; the Critic did
not edit the manuscript.  Shared rules required hash validation before every
role claim, prohibited both roles from running in one invocation, preserved
earlier artifacts, and required selectable-text vector PDFs with visual
inspection.  The complete operational prompt appears in the workflow record
cited above.

\subsection*{Schedule, iterations, and time accounting}

Research artifacts span 10--28 July 2026.  The workflow produced 18 numbered
internal manuscript iterations; the released manuscript is intentionally
unnumbered.  From 17 July onward, the scheduler started one state-aware worker
every six hours.  Each invocation selected at most one role from the persisted
phase.  Researcher iterations required at least five active hours, Critic
iterations at least one active hour, and every successful handoff ended the
invocation.

The retained late-stage log records 40.68 hours of Researcher and Critic
workflow time, including unsuccessful and repeated handoff attempts.  The final
production-and-review cycle comprised 358.38 Researcher minutes and 65.31
Critic minutes, or 7.06 hours in total.

\begin{thebibliography}{9}
\bibitem{Ghidelli}
L.~Ghidelli,
\emph{Progress on the unfair $0$--$1$ polynomials conjecture using linear
recurrences and numerical analysis}, arXiv:2209.09843 (2022),
\url{https://arxiv.org/abs/2209.09843}.

\bibitem{Hare}
K.~G.~Hare,
\emph{Computational progress on the unfair $0$--$1$ polynomial conjecture},
Experimental Mathematics, advance online publication, 15 July 2025,
\url{https://doi.org/10.1080/10586458.2025.2527786}.

\bibitem{IS}
M.~Idris and J.-M.~Sac-\'Ep\'ee,
\emph{Algorithmic aspects of Newman polynomials and their divisors},
arXiv:2601.11486 (2026),
\url{https://arxiv.org/abs/2601.11486}.

\bibitem{IS2}
M.~Idris and J.-M.~Sac-\'Ep\'ee,
\emph{On the existence of Newman and Littlewood multiples for certain integer
polynomials}, arXiv:2607.15520 (2026),
\url{https://arxiv.org/abs/2607.15520}.

\bibitem{DJS}
P.~Drungilas, J.~Jankauskas, and J.~\v{S}iurys,
\emph{On Newman and Littlewood multiples of Borwein polynomials},
arXiv:1609.07295 (2016),
\url{https://arxiv.org/abs/1609.07295}.

\bibitem{KR}
M.~Krasner and B.~Ranulac,
\emph{Sur une propri\'et\'e des polyn\^omes de la division du cercle},
Comptes Rendus de l'Acad\'emie des Sciences de Paris \textbf{204} (1937),
397--399.

\bibitem{OP}
A.~M.~Odlyzko and B.~Poonen,
\emph{Zeros of polynomials with $0,1$ coefficients},
L'Enseignement Math\'ematique \textbf{39} (1993), 317--348.
\end{thebibliography}

\end{document}
